\scanpage{37}\mspage{24}

\uline{Lemme} \deleted{\unsure{Pour|Soit} $i_0\in I$, \unsure{$K$} \unsure{désignons par} $P_{i_0}$} \inserted{Soit \deleted{\unsure{$I$}} $P\subset I$ l'ens des} $i_0\in I$ tels que pour tt ouvert $\deleted{J}K$ de $J=I_{/i_0}$, la flèche $u_{\deleted{\unsure{J'}}K,J}$ est dans $C$. \deleted{Pour qu'on ait} \inserted{Ceci dit,} soit $i_0\in\deleted{\unsure{P}}I$ tel que tout $i<i_0$ soit dans $P$. Alors $i_0\in P$. (On suppose $u_\bullet\in \unsure{C|\mathcal C}_I$, mais $I$ quelc.)

En effet, \deleted{\unsure{pour un tel cas}} soit $\deleted{J}K\subset J=I_{/i_0}$\inserted{ un ouvert}. Si \deleted{$\unsure{J'=J}$} $\deleted{\unsure{J'}}K=J$, il est trivial que $u_{\deleted{\unsure{J'}}K,J}\in C$ (c'est un iso). Si $\deleted{\unsure{J'}}K\ne J$, on a \deleted{$K\subset\unsure{J'}$}

\[
 K\subset J'\overset{\mathrm{def}}{=}\deleted{\unsure{I}}J\setminus\{i_0\}\underset{\unsure{\ne}}{\subset}J
\]

et pour prouver que $u_{K,J}\in C$, il suffit de prouver que $u_{K,J'}$, $u_{J',J}\in C$. Or $u_{\deleted{\unsure{K}}J',J\deleted{\unsure{q}}}\in C$ par hyp. $u_\bullet\in C_I$. Et $u_{\deleted{\unsure{J}}K,J'}\in C$ \unsure{par|pour} ce qui a été prouvé pour le cas de $J'$, compte tenu que \inserted{par hyp sur $i_0$,} tout $i\in J'$ (i.e. tt $i<i_0$) \deleted{\unsure{satisfait la}} \inserted{est dans $P$}.

\deleted{\unsure{Corollaire}} \uline{Corollaire}\qquad Si $I$ est artinien, $P=I$.

Car \unsure{toute} partie $P$ de $I$ ayant la propriété qu'on vient de dire ($I_{/i_0}'\subset P\Longrightarrow i_0\in P$) est égale à $I$, \deleted{\unsure{Car}} Si en effet il y avait $x_0\notin P$, il existerait $x_1<x_0$ avec $x_0\notin P$, et par récurrence on \unsure{construirait} une suite \unsure{strictmt|strictement} décroissante dans $I$, absurde.

% archive sheet number: 36

